% Part: second-order-logic
% Chapter: syntax-and-semantics
% Section: expressive-power

\documentclass[../../../include/open-logic-section]{subfiles}

\begin{document}

\olfileid{sol}{syn}{exp}
\olsection{अभिव्यक्तीचे सामर्थ्य}

\begin{explain}
द्वितीय-क्रम
चलांवरील
संख्यकीकरणामुळे
भाषेचे अभिव्यक्तीचे
सामर्थ्य प्रथम-क्रम
तर्कशास्त्राच्या
तुलनेत प्रचंड
वाढते.
द्वितीय-क्रम
अस्तित्ववाचक
संख्यकीकरणाने
विशिष्ट गुणधर्म
असलेली फलने
किंवा संबंध
अस्तित्वात आहेत,
असे म्हणता येते.
प्रथम-क्रम
तर्कशास्त्रात
हे करण्याचा
एकमेव मार्ग
म्हणजे त्यासाठी
अतार्किक चिन्ह
(म्हणजे !!a{function}
किंवा !!{predicate})
निर्दिष्ट करणे.
द्वितीय-क्रम
वैश्विक संख्यकीकरणाने
आधारक्षेत्राचे
सर्व उपसंच,
त्यावरील सर्व
संबंध किंवा
!!{domain} वरून
!!{domain} मध्ये
जाणारी सर्व
फलने यांना
एखादा गुणधर्म
आहे, असे
म्हणता येते.
प्रथम-क्रम
तर्कशास्त्रात
भाषेतील अतार्किक
चिन्हांपैकी
एखाद्याला
नेमून दिलेल्या
उपसंचांना,
संबंधांना किंवा
फलनांनाच
गुणधर्म आहे,
एवढेच म्हणता
येते.
एखादा गुणधर्म
असलेले उपसंच,
संबंध किंवा
फलने अस्तित्वात
आहेत किंवा
त्या सर्वांना
तो गुणधर्म
आहे असे म्हणताना,
त्या गुणधर्माचे
वर्णन करतानाही
द्वितीय-क्रम
संख्यकीकरण
वापरता येते.
यामुळे प्रथम-क्रम
तर्कशास्त्रात
परिभाषित न
होणारे संबंध
परिभाषित करता
येतात आणि
तेथे व्यक्त न
होणारे आधारक्षेत्राचे
गुणधर्म व्यक्त
करता येतात.
\end{explain}

\begin{defn}
$\Struct{M}$ ही
भाषा~$\Lang{L}$
साठी !!a{structure}
असेल, तर
$R \subseteq \Domain{M}^2$
हा संबंध~$\Lang{L}$
मध्ये \emph{परिभाष्य}
असतो, जर फक्त
$x$ आणि~$y$
हीच मुक्त चले
असलेले असे
!!{formula}~$!A_R(x, y)$
असते की
$s(x) = a$
आणि $s(y) = b$
असताना
$R(a, b)$
खरे असते
(म्हणजे
$\tuple{a,b} \in R$),
जर आणि फक्त
जर
$\Sat{M}{!A_R(x, y)}[s]$.
\end{defn}

\begin{ex}
प्रथम-क्रम
तर्कशास्त्रात
एकरूपता संबंध~
$\Id{\Domain{M}}$
(म्हणजे
$\Setabs{\tuple{a,a}}{a \in
  \Domain{M}}$)
सूत्र~$\eq[x][y]$
याने परिभाषित
करता येतो.
द्वितीय-क्रम
तर्कशास्त्रात
हा संबंध
\emph{$\eq$ शिवाय}
परिभाषित
करता येतो.
कारण $a$
आणि $b$
हे~$\Domain{M}$
चे तेच
!!{element}
असतील, तर
ते~$\Domain{M}$
च्या त्याच
उपसंचांचे
!!{element}s
असतात (कारण
संच त्यांतील
!!{element}s
नी ठरतात).
उलट $a$
आणि $b$
भिन्न असतील,
तर ते त्याच
उपसंचांचे
!!{element}s
नसतात:
उदा.,
$a \in \{a\}$
पण $a \neq b$
असताना
$b
\notin \{a\}$.
म्हणून
“~$\Domain{M}$
च्या त्याच
उपसंचांचे
!!{element}s
असणे” हा
संबंध $a$
आणि $b$
यांच्याबाबत
$a = b$
असेल तर आणि
तरच लागू
होतो.
हा संबंध
द्वितीय-क्रम
तर्कशास्त्रात
व्यक्त करता
येतो, कारण
~$\Domain{M}$
च्या सर्व
उपसंचांवर
संख्यकीकरण
करता येते.
म्हणून पुढील
!!{formula}
$\Id{\Domain{M}}$
परिभाषित करते:
\[
\lforall[X][(X(x) \liff X(y))]
\]
\end{ex}

\begin{prob}
$\lforall[X][(X(x) \lif X(y))]$
(लक्षात घ्या:
$\liff$ नव्हे,
$\lif$!)
हे सूत्र
$\Id{\Domain{M}}$
परिभाषित करते,
हे दाखवा.
\end{prob}

\begin{ex}
$R$ हे
द्विस्थानी
!!{predicate}
असेल, तर
$\Assign{R}{M}$
हा~$\Domain{M}$
वरील द्विस्थानी
संबंध असतो.
थोडे गोंधळात
टाकणारे असले,
तरी~$R$ हेच
चिन्ह त्या
!!{predicate}
साठी~$R$ आणि
संबंध~$\Assign{R}{M}$
साठीही वापरू.
$R$ चे
\emph{संक्रमक
संवरण}~$R^+$
हा $a$
आणि $b$
यांच्यात तेव्हाच
लागू होणारा
संबंध आहे
जेव्हा काही
$c_1$, \dots, $c_k$
साठी
$R(a,c_1)$,
$R(c_1, c_2)$,
\dots, $R(c_k,b)$
लागू होतात.
यात $k = 0$
हा प्रसंगही
येतो, म्हणजे
$R(a,b)$
लागू असेल,
तर $R^+(a,b)$
ही लागू असतो.
याचा अर्थ
$R \subseteq R^+$.
खरे तर
$R^+$ हा~$R$
समाविष्ट करणारा
सर्वांत लहान
संक्रमक संबंध
आहे.
द्वितीय-क्रम
तर्कशास्त्रात
$X$ हा~$R$
समाविष्ट करणारा
संक्रमक संबंध
आहे असे
म्हणता येते:
\begin{multline*}
  !B_R(X) \ident \lforall[x][\lforall[y][(R(x,y) \lif X(x, y))]] \land {}\\
\lforall[x][\lforall[y][\lforall[z][((X(x,y) \land X(y,z)) \lif X(x,
      z))]]].
\end{multline*}
पहिला संयुक्तांश
$R \subseteq X$
सांगतो आणि
दुसरा $X$
संक्रमक आहे
असे सांगतो.

$X$ हा
असा सर्वांत
लहान संबंध
आहे असे
म्हणणे म्हणजे
$R$ समाविष्ट
करणाऱ्या
आणि संक्रमक
असलेल्या
प्रत्येक संबंधात
तो समाविष्ट
आहे असे
म्हणणे.
म्हणून
पुढील
!!{formula}
याने~$R$
चे संक्रमक
संवरण
परिभाषित
करता येते:
\[
R^+(X) \ident !B_R(X) \land \lforall[Y][(!B_R(Y) \lif
  \lforall[x][\lforall[y][(X(x, y) \lif Y(x,y))]])].
\]
$\Sat{M}{R^+(X)}[s]$
हे
$s(X) = R^+$
असेल तर आणि
तरच खरे.
$R$ चे
संक्रमक संवरण
प्रथम-क्रम
तर्कशास्त्रात
व्यक्त करता
येत नाही.
\end{ex}

\end{document}
